\section{MDL 把拟合与描述长度放在同一尺度}
\label{sec:13-Machine-Learning/11-Information-Theoretic-Learning/03}

你手头有一组带噪声的数据点，想选一个多项式去解释它。一次多项式只抓到趋势，残差大；九次多项式穿过几乎所有训练点，残差几乎为零。但你并不觉得九次多项式更好——它只是在死记硬背每个点的位置，没有发现可复用的规律。

问题是：你怎么把这种直觉变成两个模型之间公平的数字比较？拟合误差有自己的单位（均方误差、负对数似然），复杂度的单位又该是什么？参数个数？阶数？自由参数的计数不是同一个量纲，没法直接加在误差上。

\subsection{把模型和残差都拉进同一根编码尺}

Minimum description length（MDL）的做法很直接：把选模型变成一场通信游戏。你想把数据告诉接收者，但你们事先约定了一个编码协议。如果你可以先传一个模型，再用这个模型帮接收者预测数据，接收者需要的数据补充就会变短。总传输量是最公平的比较尺度——因为无论模型复杂度还是拟合残差，最后都化成了同一个东西：比特数。

最简单的 two-part MDL 选择最小化
\[
L(M) + L(D\mid M)
\]
的模型。第一项编码模型的结构和参数，第二项在已知模型后编码数据或残差。复杂模型的第二项短（残差小），但第一项长（要把模型自身交代清楚）；简单模型反过来。总码长最低的那个解释，就是在信息论意义上最经济的解释。

这个框架把``好模型''翻译成了``好的压缩工具''：抓住了规律的模型能帮接收者省码长；只靠死记的模型把自己的庞大参数传过去之后，接收者并没有获得任何压缩数据的杠杆。

\subsection{Likelihood 怎么变成码长}

概率模型和编码之间有一个干净的操作性对应：给事件 \(D\) 分派概率 \(p(D\mid M)\) 的模型，可以给 \(D\) 分配一个理想码长
\[
L(D\mid M) \approx -\log_2 p(D\mid M).
\]
底数 2 表示以 bit 为单位的码长（第一节已经说明：log 底只改变数值单位，不影响比较结果）。`理想'和 \(\approx\) 是因为真实码长必须是整数——逐符号编码与 \(-\log_2 p\) 相差不到 1 bit；对长序列，算术编码可以在平均意义上逼近它。

把数据看成从某个真实源头产生的，再取期望，这正是第一节``用模型分布编码真实来源''的平均码长——即交叉熵。

高 likelihood 意味着模型把短码分配给了实际出现的数据。但仅凭这一点还不够。一个拥有大量自由参数的模型可以通过调参把任何数据都赋以高 likelihood——接收者如果不知道这些参数的具体取值，他无法单靠残差码长重建设数据。第一项 \(L(M)\) 的存在，正是为了防止你选一个``对什么数据都能说出高概率''的巨无霸模型。

\subsection{prefix code 和 Kraft 不等式：概率和码长的双向可译}

码长得满足可译条件，否则接收者拿到一串 bit 不知道怎么切分。prefix code（前缀码）——任何码字都不是另一个码字的前缀——正是满足这一点的最自然构造。prefix code 的码长必须遵守 Kraft 不等式：
\[
\sum_x 2^{-L(x)} \le 1.
\]
这个不等式的来源是前缀码的树结构：每个码字占据树上的一个叶子，编码树从根到叶子的路径被独占。两个方向的严格讨论见\hyperref[sec:07-Information-Theory/04-Source-Coding/01]{``前缀码与 Kraft 不等式''}。

Kraft 不等式真正的力量在于它把码长和概率挂在了一起。从码长出发，\(2^{-L(x)}\) 可以视为一个 sub-probability（和不超 1 的非负权重）。从概率出发，任何概率分布 \(p(x)\) 都可以诱导一个码长接近 \(-\log_2 p(x)\) 的可译码。于是 MDL 的三个层面——概率建模、编码可实现性和模型比较——被同一套数学锁紧了：其中的 bit 不是比喻，而是真实的编码操作。

\subsection{连续参数的编码不是参数个数加个 log}

假设两个模型都有三个参数。直觉告诉你编码成本相同。但仔细看：

第一，连续参数需要指定精度。把参数 \(\theta\) 量化到有限精度 \(\Delta\)，码长大约与 \(-\log_2\Delta\) 成正比。精度要求越高，码长越长。而不同模型对参数精度的敏感度不同——有些模型中参数稍微偏差数据 likelihood 就暴跌，必须高精度编码；有些模型中参数变化很大时 likelihood 只微降，低精度就够。参数个数相同，编码成本可以差几个数量级。

第二，参数的可识别性。如果两个参数交换后模型完全相同（如神经网络的权重对称性），逐一编码会重复计费——同一个等价函数被编码了多次。

第三，BIC 可以视为 MDL 在大样本、正则模型、固定维度下的一种近似。但奇异模型（如混合模型，Fisher 信息矩阵退化）、深度网络（大量对称和参数冗余）以及有限样本下，BIC 的近似可能失准。Normalized maximum likelihood（NML）、Bayesian mixture code 和 prequential code 各自对``复杂度''有不同的操作化定义与有限样本行为，不能混为一谈。

prequential 编码的直觉尤其好：先用前一段数据训练模型，再用模型给下一段数据编码，持续滚动更新。一个能快速学到规律的模型，会较早地给未来数据分配短码；一个只能记住过去数据的模型，在每一个新段到来时都要付出高预测代价。此时训练算法本身也成为编码方案的一部分——模型是否有用，由它在未见数据上的压缩能力说了算。

\subsection{在真实数据上拨一遍算盘}

把刚才的框架放到一个具体的例子中：用多项式解释温度与传感器漂移的关系。

阶数从 1 往上提。训练残差的负对数似然一路下降，这在意料之中——多项式越灵活，穿过数据的能力越强。但同时，你需要编码阶数本身（一个整数）、每个系数的具体值（量化到某精度），以及残差分布的参数（如 Gaussian 的 \(\sigma\)）。

低阶模型：结构码长很短（三个系数 + 阶数），但残差大，数据码长很长。高阶模型：数据码长骤降，但参数码长攀升——你要传十几个系数和它们的高精度值。两条曲线的交叉点通常出现在中间某阶，那是最小总码长对应的模型。

但这个最低点不是绝对的。它依赖两个关键选择：

一是残差的噪声模型。如果残差有重尾但用了 Gaussian 编码，少数异常点会被赋以极长的码——Gaussian 的尾巴薄，罕见事件的 surprise 极高。换成重尾编码（如 Laplace 或 t 分布），那些异常点的码长立刻缩短，可能比盲目提高多项式阶数更划算。MDL 要求你把``模型结构''和``解释不了的剩余''都说清楚，但它不会自动替你选哪个模型族是合理的——选了不合适的噪声模型，总码长就会偏好不必要的复杂结构。

二是参数的量化精度。精度翻倍，每个参数的码长增加 1 bit。在某些精度下最低点落在三阶，换了精度可能挪到四阶。这不是说 MDL 不稳定，而是说``用哪种编码''本身就是建模的一部分。

实际比较时的好习惯：固定数据精度、log 底数和编码协议，分别报告结构、参数和数据三部分的码长，让读者看到每一项的贡献。用 held-out 或 prequential code 做敏感性检查，确认改变参数量化精度后结论方向不变。

\subsection{同一条拟合曲线，不同的账单会选出不同的模型}

把刚才的讨论收束为一个可操作的观点。对同一组数据，随多项式阶数升高：

\begin{itemize}
\tightlist
\item 训练 NLL 几乎总是单调下降——灵活模型给数据更高似然；
\item 参数码长几乎总是单调上升——更多自由度要更多的 bit 交代；
\item 总码长先降后升，在某个中间阶数取最低点。
\end{itemize}

但``训练 NLL 最小''选出的模型、``BIC 最小''选出的模型和某个特定 two-part code 总码长最小选出的模型，不必是同一个模型。它们使用的是不同账单——训练 NLL 根本没计参数码长，BIC 用渐进公式近似了参数码长，two-part code 则需要显式指定量化方案和噪声模型。这并非同一原则的数值误差，而是三套不同的计费标准。MDL 的力量不在于给出一个不可辩驳的单一答案，而在于把所有隐性的选择——编码方案、噪声分布、精度——全部拉到台面上，让比较变得透明。
